% SGA 3, Expose VI_B, Loop 1 English reconstruction.
% Scope: Corollary 7.1.1; Definitions and Remarks 7.2(i)--(v);
% and Corollary 7.2.1 with their proofs, including source
% editor's notes 70--72.
% Authority: Polo--Gille Exp6B-13oct24.pdf, component-local pages
% 51--52 (combined-reader pages 392--393; re-edition headers
% 381--382; original-margin printed pages 363--365).
% Definition 7.2.2 begins on local page 52 and is excluded.
% OCR and Reinhold English Markdown are comparison witnesses only.
% This unit contains no diagrams.

\medskip
\phantomsection\label{cor:sga3-VIB-7.1.1}%
\noindent\textbf{推论 7.1.1.---}%
\footnote{\label{note:sga3-VIB-70}编者注: 增补本推论
以指出 \SGARefLink{prop:sga3-VIB-7.1}{7.1} 的这一特殊情形.}
\emph{设 \(G\) 为局部有限型 \(k\)-群,
\(A\) 和 \(B\) 为 \(G\) 的两个平滑
有限型 \(k\)-子群, \(i_A\) (相应地, \(i_B\)) 为
\(A\) (相应地, \(B\)) 到 \(G\) 的包含态射. 假设
\(B\) 正规化 \(A\). 则}
\phantomsection\label{form:sga3-VIB-7.1.1-generated-product}%
\[
  A\cdot B=\Gamma_G(i_A,i_B).
\]

\emph{证明.}
设
\[
  H=A\rtimes B
\]
为 \(A\) 和 \(B\) 的半直积 (参见 I, \SGARefLink{sga3:i:definition:2-3-5}{2.3.5});
它是平滑有限型 \(k\)-群. 因此群态射
\[
  u:H\longrightarrow G,
  \qquad
  (a,b)\longmapsto ab,
\]
拟紧. 所以由
\hyperref[prop:sga3-VIB-1.2]{1.2},
\[
  u(H)=A\cdot B
\]
为 \(G\) 的闭约化子概形, 且为
范畴 \((\mathrm{Sch}/k)_{\mathrm{red}}\) 中的群. 由
EGA~IV\(_3\), 11.10.7, \(A\cdot B=u(H)\) 几何
约化, 因为 \(H\) 几何约化; 故它是 \(G\) 的闭子群.
显然的包含关系
\[
  A\cdot B\subset\Gamma_G(i_A,i_B)
\]
即证明了本推论.

\par
\ifdim\dimexpr\pagegoal-\pagetotal\relax<12\baselineskip
  \newpage
\fi
\medskip
\phantomsection\label{defrem:sga3-VIB-7.2}%
\noindent\textbf{定义及\SGARefLink{defrem:sga3-VIB-7.2}{注 7.2}.---}%
\footnote{\label{note:sga3-VIB-71}编者注: 将原文的
\textup{(viii)} 项移至
\textup{(i\(^{\prime}\))}, 并将 \textup{(vi)} 和
\textup{(vii)} 项单列为下面的推论~\hyperref[cor:sga3-VIB-7.2.1]{7.2.1}
及定义~\hyperref[def:sga3-VIB-7.2.2]{7.2.2}.}
\begin{enumerate}
\item[\textup{(i)}]\phantomsection\label{def:sga3-VIB-7.2-i}
  设 \(G\) 为 \(k\)-群, \((X_i)_{i\in I}\) 为
  一族几何约化的 \(k\)-概形, 并对每个
  \(i\in I\), 设 \(f_i:X_i\to G\) 为 \(k\)-态射.
  \emph{由族 \((f_i)_{i\in I}\) 生成的 \(G\) 的闭子群概形}
  是 \(G\) 的最小闭子群概形,
  使得每个 \(f_i\) 都经由它分解; 在本小节中
  将其记为
  \[
    \Gamma_G\bigl((f_i)_{i\in I}\bigr).
  \]
  若 \(X\) 为 \(G\) 的子概形, 且在
  \(k\) 上几何约化, \(f:X\hookrightarrow G\) 为浸入, 则用
  \(\Gamma_G(X)\) 代替 \(\Gamma_G(f)\).

\item[\textup{(i\(^{\prime}\))}]
  \phantomsection\label{def:sga3-VIB-7.2-iprime}
  采用 \hyperref[concl:sga3-VIB-7.1-ii]{7.1\textup{(ii)}} 的记号,
  有时令
  \phantomsection\label{form:sga3-VIB-7.2-iprime-Gamma}%
  \[
    \Gamma'_G(f)=\bigcup_{n\geq1}f^n(X^n).
  \]
  注意 \(\Gamma'_G(f)\) 是 \(G\) 的一个子集,
  它在 \hyperref[par:sga3-VIB-3.0]{3.0} 的意义下关于群运算稳定.

\item[\textup{(ii)}]\phantomsection\label{rem:sga3-VIB-7.2-ii}
  若 \(X_1\) 和 \(X_2\) 为几何约化的
  \(k\)-概形, 且 \(f_1:X_1\to G\) 和 \(f_2:X_2\to G\) 为
  \(k\)-态射, 使得在集合意义下
  \[
    f_1(X_1)=f_2(X_2)
  \]
  成立, 则显然
  \[
    \Gamma_G(f_1)=\Gamma_G(f_2).
  \]

\item[\textup{(iii)}]\phantomsection\label{def:sga3-VIB-7.2-iii}
  设 \(E\) 为 \(G\) 的一个子集, 使得 \(G\) 的约化
  子概形 \(\overline E\) 几何约化.
  \emph{由 \(E\) 生成的 \(G\) 的闭子群概形}
  是子群概形
  \[
    \Gamma_G(E):=\Gamma_G(i),
  \]
  其中 \(i:\overline E\hookrightarrow G\) 为
  该约化子概形的包含态射.

\item[\textup{(iv)}]\phantomsection\label{rem:sga3-VIB-7.2-iv}
  由于 \(G\) 的每个子群概形都闭, 见
  \hyperref[sga3:VIA:0.5.2]{VI\(_A\), 0.5.2},%
  \footnote{\label{note:sga3-VIB-72}编者注: 这里删去了
  \(G\) 在 \(k\) 上局部有限型的
  假设.}
  我们将简称 ``生成的子群概形'',
  而不说 ``生成的闭子群概形''.

\item[\textup{(v)}]\phantomsection\label{rem:sga3-VIB-7.2-v}
  设 \(X\) 为几何约化的 \(k\)-概形,
  \(f:X\to G\) 为 \(k\)-态射. 假设 \(f(X)\)
  包含 \(G\) 的单位元 \(e\). 令
  \[
    X'^{\,1}=X\times_kX,
    \qquad
    f'^{\,1}=\mu\circ\bigl(f\times_k(c\circ f)\bigr),
  \]
  并对 \(n>1\), 令
  \phantomsection\label{form:sga3-VIB-7.2-v-recursion}%
  \[
    X'^{\,n}=X'^{\,1}\times_kX'^{\,n-1},
    \qquad
    f'^{\,n}
      =\mu\circ\bigl(f'^{\,1}\times_k f'^{\,n-1}\bigr).
  \]
  则 \(\Gamma_G(f)\) 为 \(G\) 的约化子概形,
  其底空间为
  \[
    \bigcup_{n\geq1}f'^{\,n}(X'^{\,n}).
  \]
  的闭包.

  事实上, 回顾
  \hyperref[concl:sga3-VIB-7.1-ii]{7.1\textup{(ii)}} 的记号:
  \[
    X^1=X\amalg X,
    \qquad
    X^n=X^1\times_kX^{n-1}\quad(n\geq2).
  \]
  由假设 \(e\in f(X)\), 对每个 \(n\geq1\), 有
  \phantomsection\label{form:sga3-VIB-7.2-v-inclusions}%
  \[
    f^n(X^n)
      \subset f'^{\,n}(X'^{\,n})
      \subset f^{2n}(X^{2n}).
  \]
  因此, 存在一个整数 \(n\) 使得
  \(f^n(X^n)=\Gamma_G(f)\) 当且仅当存在一个
  整数 \(m\) 使得
  \(f'^{\,m}(X'^{\,m})=\Gamma_G(f)\).
\end{enumerate}

由注~\hyperref[rem:sga3-VIB-7.2-v]{7.2\textup{(v)}}
得到下述结果.

\medskip
\phantomsection\label{cor:sga3-VIB-7.2.1}%
\noindent\textbf{推论 7.2.1.---}
\emph{设 \(X\) 为几何约化且几何
连通的 \(k\)-概形, \(f:X\to G\) 为 \(k\)-态射,
使得 \(f(X)\) 包含 \(G\) 的单位元. 则
\(k\)-群 \(\Gamma_G(f)\) 连通, 因而
不可约.}

\emph{证明.}
每个 \(X'^{\,n}\) 均连通, 所以所有
\[
  f'^{\,n}(X'^{\,n}),
\]
的并连通, 因为它们都包含单位元. 因此其闭包
\(\Gamma_G(f)\) 也连通. 当 \(G\) 从而
\(\Gamma_G(f)\) 在 \(k\) 上局部有限型时,
由 \hyperref[prop:sga3-VIA-2.4]{VI\(_A\), 2.4},
而在一般情形中由
\hyperref[thm:sga3-VIA-2.6.5]{VI\(_A\), 2.6.5\textup{(iii)}},
可知 \(\Gamma_G(f)\) 不可约.

% 翻译统计: 源行数 193; 输出行数 194 (含本统计行); 行内/行间数学片段 108; 翻译占位 0. 底本疑点另报主线程.
